\scanpage{22}
\mspage{\unsure{11|II}}
Donc l'image de $A$ est $P$. Reste à vérifier
que le noyau de \deleted{$A$ est} $u:A\to P$ \deleted{\unsure{est|$\to$}} [qui, on
le sait, est fini sur $S$, fibre générique étant
le groupe unité] est le groupe unité sur $S$.
En tous cas \deleted{ce groupe est d'ordre} \deleted{\unsure{la fib|Ce fib.}} si
$p$ est l'exposant caractéristique du corps résiduel
$k$, alors la fibre $K_0$ est un $p$-groupe.

\deleted{En effet, si $\xi_0$ est un pt de $K_0$ \unsure{rat./k}
[N.B. on peut supposer $k$ alg. clos \inserted{et $V$ complet}], d'ordre
$n$ premier à $p$, alors $\xi_0$ provient
d'une section $\unsure{\xi|\zeta}$ d'ordre $n$ de $A$ sur $S$,
qui donne une section d'ordre $n$ de
\unsure{le|$\xi$|\&}
\diagram{\begin{tikzcd}[ampersand replacement=\&,column sep=small,row sep=small]
\deleted{Y}\arrow[rr,bend left=13,dash]\arrow[dr] \& \& \deleted{X}\arrow[dl]\\
\&\deleted{S}\&
\end{tikzcd}}
$\uline{\mathrm{Pic}}_{X/S}$ nulle au pt fermé $s$, \unsure{Donc|donc}}

Cela résulte facilement du fait que si
$n$ est premier à $p$, alors $(\uline{\mathrm{Pic}}_{X/S})_n$ est
\deleted{\uline{étale}} \inserted{\uline{non ramifié}} sur $S$, et que \deleted{\unsure{$A$ est|$A_n$}} $A_n$ est \deleted{\unsure{ét|ét.}}
un \uline{revêtement étale} de $S$. Donc si $p=1$
i.e. $k(\unsure{v|V})$ est de car $0$, on a fini.
Sinon, il faut regarder de plus près \unsure{,|.|\, } - - - -

\margin{?}À étudier: noyaux, images d'homomorphismes de
schémas abéliens et multiabéliens, en car $\deleted{\unsure{k|\ne}} \unsure{\ne|k} \unsure{0}$,
notamment sur une base non réduite, ou
% archive sheet number: 21
